\chapter{Proof that a Symmetric Energy Distribution fails to cancel in the Product-Sum}
\chaptermark{Product-Sum Proof}
\label{ch:symm_sample}
	In the ansatz reflected in equation (\ref{eq:simple_sym_eec}), a symmetric set with a distribution roughly approximating a gaussian was utilized, however, we can go beyond that to properly show this result is the case in a set drawn from a generalized Gaussian distribution
	\newline
	\begin{proof}{}
		We begin with a gaussian distribution $P(x)$ and sample points from it such that we construct a "Gaussian sample set" with the following properties 
		\newline
		\begin{definition}[Gaussian Sample Set]
		\label{def:GSS}
		Take a set S that is constructed of points randomly sampled such that  
		\begin{equation}
			S \approx P(x)=  e^\frac{-x^2}{\sigma}
		\end{equation}
		Which gives the following properties
		\begin{enumerate}
		\item{ $ S= \left\{ x | \bar{x} = \delta_{\bar{x}} \right\}$} 
		\item{ $ \frac{1}{N} \sum_{x\in S} x^2 = \left( \sigma + \delta_\sigma \right)^2$}
		\item{ $ \left| \delta_{\bar{x}} \right|, \left| \delta_\sigma \right| << 1 $}
		\item{ $\lim_{N\rightarrow \inf} \left| \delta_{\bar{x}} \right|, \left| \delta_\sigma \right| = 0 $}
		\end{enumerate}
		
		\end{definition}
		
		Then, the Product-Sum is calculated by following the standard algorithm 
		\newline
		\begin{algo}[Product-Sum]
		\begin{equation}
			\begin{split}
				\forall{x} \in {S} & s.t. len(S) > 1\\
				& S = S / x  \\
				& PS = PS + ( x * y ) \hspace{1.5 cm} \forall y \in S \\	
			\end{split}	
		\end{equation}
		\textit{Product-sum algorithm for the sample set S. Initially the PS varaiable is set at 0}
		\end{algo}
		Using the properties of S noted in definition \ref{def:GSS}, one notes that for any sampled value x, there should be approximately $n= \floor{\frac{1}{\sqrt{2 \pi (\sigma+ \delta_\sigma}} e^{\frac{\left( x - \delta_{\bar{x}} \right)^2}{\sigma+\delta_\sigma}}}$ of the same value. 
		In fact, by including these small corrections and an additional small correction $\left|\delta_{n, x}\right| << 1$  we fix this number to the Gaussian curve within some fluctations. 
	With this in mind, one breaks the contributions to the Product-sum into two categories, the $y=x$ and $y \neq x$ case. 
	For the first case, the contribution is simply 
	\begin{equation}
		PS_{y=x} = x^2 \left( n(x) - 1 \right)! 
	\end{equation}
	and in the second case 
	\begin{equation}
		PS_{y\neq x} = x y n(x) n(y)
	\end{equation}
	Thus the product-sum becomes
	\begin{equation}
		PS = \sum_{x\in S}\left[  x^2 \left(n(x) - 1\right)! + \sum_{y > x, y\in S}x y n(x)n(y) \right] 
	\end{equation}
	With a sufficiently high number of samples, we can shift to the continuum limit and write
	\begin{equation}
		\tilde{PS} = \int_{-\infty}^{\infty} \dd x \hspace{0.2 cm} \left[ x^2 \Gamma(\left(n(x) - 1\right)) + \int_{x}^{\infty} \dd y \frac{xy}{2 \pi \left( \sigma + \delta \sigma \right)} e^{\frac{- \left( x - \delta_{\bar{x}} \right)^2 - \left( y - \delta_{\bar{x}} \right)^2}{2 \left( \sigma + \delta \sigma\right)}}   \right]
	\end{equation}

	\end{proof}
